\scanpage{32}

On devine que l'iso

\[
{E^{\deleted{\unsure{\vee}}}}^{\otimes_A2}\otimes_{\unsure{k|h}}L^\vee\simeq {E^\vee}^{\otimes_A2}\otimes_k L
\]

i.e.

\[
{E^\vee}^{\otimes_A2}\simeq E^{\otimes_A2}\otimes_k{L^{\deleted{\unsure{\vee}}}}^{\otimes_k(-2)}
\]

est déduit de l'iso can. (qui respecte les structures de $A$-modules)

\[
E^\vee\simeq E\otimes_k L^{\otimes_k(-1)}
\]

en \unsure{prenant|élevant} les \deleted{\unsure{puis|par}} carrés tensoriels sur $A$.

L'iso de $A$-modules cherché (mod scalaires provenant de $\underline O_{\unsure{\alpha|X}}^*$)

\[
E\simeq\underline{\operatorname{End}}(E)/\mathcal A\qquad(\simeq E^{\otimes_A2}\otimes_k L^\vee)
\]

équivaut à la \deleted{\unsure{iso}}donnée d'une section de

\[
E\otimes_k L^\vee,
\]

base de ce $\mathcal A$-module mod section provenant de $O^*$, Or il n'y en a pas de canonique, visiblement, \deleted{Donc} \inserted{i.e.} il n'y a pas de sous-fibré en droites canonique $D\subset E$ associée à la structure de similitude de départ sur $E$. Donc la supposition du début est erronée -- $\operatorname{End}(E)/\mathcal A$ est, modulo \deleted{\unsure{m}}Module inversible provenant de la base, un \deleted{\unsure{m}}Module inversible \uline{carré}, et c'est la racine carrée qu'il faudrait arriver à saisir \ldots
% archive sheet number: 31
